\section{Conclusion}
\label{sec:conclusion}

We presented UCDPA (Uncertainty-Calibrated Dual-Path Adaptation), a test-time adaptation method that addresses three limitations of the COME baseline: uniform treatment of samples, no pseudo-label leverage, and no mechanism linking softmax and evidential uncertainty. UCDPA introduces (i) the UCB (Uncertainty-Confidence Balance) score, a differentiable per-sample routing signal that is monotonically increasing in softmax confidence and monotonically decreasing in Dirichlet uncertainty (Theorem~\ref{thm:ucb}); (ii) a dual-path adaptation loss that softly routes each sample between conservative opinion-entropy minimization (for uncertain samples) and pseudo-label cross-entropy (for confident samples), with class-balanced reweighting; (iii) a calibration regularizer that maintains batch-mean consistency between softmax confidence and evidential uncertainty (Proposition~\ref{prop:cal}); and (iv) a feature prototype alignment term that stabilizes the representation under shift.

On ImageNet-C with ViT-B/16, UCDPA improves mean corruption accuracy by \ucdpadelta{} percentage points over COME (\ucdpaacc{} vs.\ \comeacc{}, $p < \pvalue$), with improved calibration (ECE \ucdpaece{} vs.\ \comeece{}), improved uncertainty quantification (MI \ucdpami{} vs.\ \comemi{}), and improved error detection (AUROC \ucdpaauroc{} vs.\ \comeauroc{}). UCDPA also outperforms recent baselines including CoTTA, MEMO, POEM, and FOA (Table~\ref{tab:cross-backbone}). Ablations show that UCB routing and class-balanced reweighting provide the largest gains; the calibration regularizer contributes marginally at this scale but serves as a theoretically-motivated safeguard against drift under severe distribution shift, and FPA is an optional auxiliary (disabled by default for efficiency).

\paragraph{Future work.}
Three directions seem most promising. First, \textbf{complete cross-backbone validation}: while we have validated UCDPA on ViT-B/16 (LayerNorm) and CIFAR-ResNet-18 (BatchNorm), a full ResNet-50/ImageNet-C evaluation with a valid source checkpoint is needed to close the cross-architecture gap. Second, \textbf{online threshold adaptation}: replacing the fixed UCB thresholds with online quantile-based or bandit-based adaptation, which could improve robustness to varying shift severity without introducing instability. Third, \textbf{conformal prediction}: leveraging the calibrated uncertainty for distribution-free selective prediction at test time, providing formal coverage guarantees on the abstention set---a natural extension of the error-detection results (Table~\ref{tab:auroc}) to a setting with guarantees. We also see potential in extending the calibration regularizer to a per-bin (rather than batch-mean) form, which could provide tighter calibration at the cost of additional computation.
